% ============ 02. Contexto y guía de lectura ============
\section{Contexto y guía de lectura}

\subsection{Qué es este documento}

Este documento desarrolla el \textbf{criterio 5 de la rúbrica}:
\emph{Lógica de creación del repositorio y cumplimiento del objetivo principal}. El repositorio auditado es el
Observatorio de investigadores reconocidos del Ministerio de Ciencia,
Tecnología e Innovación de Colombia (MinCiencias), que consolida y analiza
\textbf{6 convocatorias (2013--2021)} del padrón oficial de investigadores
y su producción científica. Esta serie cubre los \textbf{6 criterios} de la
rúbrica en documentos separados, más el documento maestro que los integra.

\textbf{En resumen:} Este criterio evalúa cómo se construyó el repositorio y si la ejecución fue correcta. Conclusión: la arquitectura y la lógica general son correctas y los entregables cumplen el objetivo; la ejecución fue buena en lo analítico pero no óptima en la industrialización (automatización, despliegue y pruebas), con varios flujos concretos que conviene corregir.


\begin{tcolorbox}[nota]\textbf{Versión LITE:} este documento resume el criterio (veredicto y hallazgos de mayor impacto). El desarrollo completo está en la versión completa del criterio y en el documento maestro.\end{tcolorbox}

\begin{tcolorbox}[nota]
\textbf{Cómo leer este documento:} cada PDF de esta serie desarrolla \emph{un}
criterio de la rúbrica. Los demás criterios se tratan a fondo en sus propios
documentos y en el \textbf{documento maestro} (\texttt{maestro\_auditoria.pdf}).
Si este es tu primer acercamiento, lee primero la portada y este contexto; al
final encontrarás las conclusiones del criterio.
\end{tcolorbox}

\subsection{Glosario (solo los términos que usa este documento)}

Esta tabla incluye únicamente los términos técnicos que aparecen en este
documento, explicados en lenguaje sencillo:

\begingroup
\small
\setlength{\tabcolsep}{3pt}
\begin{longtable}{>{\raggedright\arraybackslash}p{4.3cm}>{\raggedright\arraybackslash}p{9.8cm}}
\caption{Glosario de términos presentes en este documento}
\label{tab:glosario}\\
\toprule
\textbf{Término} & \textbf{Qué significa} \\
\midrule
\endfirsthead
\toprule
\textbf{Término} & \textbf{Qué significa} \\
\midrule
\endhead
    CI/CD y Docker & Automatización de pruebas y despliegue (CI/CD) y empaquetado de la aplicación en un contenedor (Docker) para que funcione igual en cualquier máquina.
 \\
\bottomrule
\end{longtable}
\endgroup
